\documentclass[12pt,reqno]{amsart}

\usepackage[T1]{fontenc}
\usepackage{lmodern}
\usepackage{microtype}
\usepackage{mathtools}
\usepackage{amssymb}
\usepackage{booktabs}
\usepackage{enumitem}
\usepackage{xcolor}
\usepackage[colorlinks=true,
  linkcolor=blue!55!black,
  citecolor=green!40!black,
  urlcolor=blue!60!black]{hyperref}

\newcommand{\Q}{\mathbf{Q}}
\newcommand{\C}{\mathbf{C}}
\newcommand{\M}{\mathcal{M}}
\newcommand{\Mbar}{\overline{\mathcal{M}}}
\newcommand{\BM}{\mathrm{BM}}
\newcommand{\RH}{\mathrm{RH}}
\newcommand{\Mod}{\operatorname{Mod}}
\newcommand{\PMod}{\operatorname{PMod}}
\newcommand{\Aut}{\operatorname{Aut}}
\newcommand{\vcd}{\operatorname{vcd}}
\newcommand{\im}{\operatorname{im}}
\newcommand{\Gr}{\operatorname{Gr}}

\newtheorem{theorem}{Theorem}[section]
\newtheorem{proposition}[theorem]{Proposition}
\newtheorem{lemma}[theorem]{Lemma}
\newtheorem{corollary}[theorem]{Corollary}
\theoremstyle{definition}
\newtheorem{definition}[theorem]{Definition}
\theoremstyle{remark}
\newtheorem{remark}[theorem]{Remark}

\title[Degree-sixteen cohomology of $\Mbar_{6,6}$]
  {The degree-sixteen cohomology of $\overline{\mathcal M}_{6,6}$ is Tate}
\author[GPT-5.6 et al.]
  {GPT-5.6, Claude Opus 5, Grok 4.6\\
   Qwen 3.8, and Daisuke Ken Sakai}
\date{1 September 2026}

\begin{document}

\begin{abstract}
We determine the Hodge type of the degree-sixteen rational cohomology of the
moduli stack of stable genus-six curves with six ordered markings.  We prove
that
\[
  H^{16}(\Mbar_{6,6};\Q)\cong\Q(-8)^{\oplus b_{16}}
\]
as a polarizable rational Hodge structure.  The argument combines the
boundary exact sequence of Canning--Larson--Payne with a calculation of its
open and boundary terms.  For the open term, their description of the
primitive quotient of lowest-weight cohomology, together with Harer's virtual
cohomological dimension, reduces the calculation to the inclusive endpoint
of the known range for the Chow--K\"unneth generation property.  For the
boundary term, we
prove by induction that $H^{16}(\Mbar_{g,n};\Q)$ is Tate for
$g\leq5$, $2g+n\leq18$, and also for $g=6$, $n\leq5$.  This includes the
nonseparating boundary atlas $\Mbar_{5,8}$; all separating boundary atlases
are handled by the K\"unneth formula, the vanishing of odd cohomology through
degree nine, and the known Tate result in even degrees at most fourteen.
The proof determines Hodge type but does not compute $b_{16}$ and does not
assert that the compact cohomology is generated by algebraic or tautological
classes.
\end{abstract}

\maketitle

\section{Introduction}

Let $\Mbar_{g,n}$ denote the smooth proper Deligne--Mumford stack of stable
complex curves of genus $g$ with $n$ ordered marked points, and let
$\M_{g,n}\subset\Mbar_{g,n}$ be its open substack of smooth pointed curves.
All cohomology and homology groups in this paper have rational coefficients.
They are considered together with their natural mixed Hodge structures.

The rational cohomology of $\Mbar_{g,n}$ is known to be Tate in every even
degree at most fourteen, uniformly in $g$ and $n$
\cite[Theorem~1.5(3) and the discussion following it]{CLP}.  In a recent
survey, Payne records that the corresponding general predictions in degrees
$16$, $18$, and $20$ remain open \cite[Section~4]{PayneSurvey}.  The purpose
of this paper is to establish a degree-sixteen case at the first of these
open degrees.

Our main theorem is the following.  The integer $b_{16}$ in its statement is
the sixteenth Betti number; no formula for it is asserted.

\begin{theorem}\label{thm:main}
Let
\[
  b_{16}:=\dim_{\Q}H^{16}(\Mbar_{6,6};\Q).
\]
There is an isomorphism of polarizable rational Hodge structures
\[
  H^{16}(\Mbar_{6,6};\Q)
  \cong \Q(-8)^{\oplus b_{16}}.
\]
\end{theorem}

The proof has two parts, corresponding to the two terms surrounding
$H_{16}(\Mbar_{6,6})$ in the boundary exact sequence.  The open part is a
lowest-weight calculation on $\M_{6,6}$.  In fact, we obtain the following
stronger description, in which $\RH^*$ denotes tautological cohomology.

\begin{theorem}\label{thm:open-intro}
For some nonnegative integer $r$,
\begin{align*}
 W_{26}H^{26}(\M_{6,6};\Q)
   &=\RH^{26}(\M_{6,6};\Q)
     \cong\Q(-13)^{\oplus r},\\
 W_{-16}H^{\BM}_{16}(\M_{6,6};\Q)
   &\cong\Q(8)^{\oplus r}.
\end{align*}
\end{theorem}

The boundary of $\Mbar_{6,6}$ contains one nonseparating component whose
ordered atlas is $\Mbar_{5,8}$, as well as separating components whose atlases
are products of smaller moduli stacks.  To control them without presupposing
the desired $(6,6)$ result, we prove the following auxiliary theorem.  A pair
$(g,n)$ is called stable if $2g-2+n>0$.

\begin{theorem}\label{thm:range-intro}
Suppose that $(g,n)$ is stable and satisfies
\begin{equation}\label{eq:range-intro}
  g\leq5\ \text{ and }\ 2g+n\leq18,
  \qquad\text{or}\qquad
  g=6\ \text{ and }\ n\leq5.
\end{equation}
Then, for $b_{g,n}=\dim_{\Q}H^{16}(\Mbar_{g,n};\Q)$,
\[
  H^{16}(\Mbar_{g,n};\Q)
  \cong\Q(-8)^{\oplus b_{g,n}}
\]
as a polarizable rational Hodge structure.
\end{theorem}

The open input for Theorem~\ref{thm:range-intro} is already supplied by the
published Chow--K\"unneth generation results except at $(g,n)=(3,12)$ and
$(5,8)$.  These two pairs and $(6,6)$ share the same numerical feature.  If
$k=2\dim\M_{g,n}-16$, then
\[
  k=\vcd(\M_{g,n})
  \quad\text{and}\quad n-1=c(g),
\]
where $c(g)$ is the marking bound in Table~1 of
Canning--Larson--Payne \cite{CLP}.  Their primitive-quotient formula peels
off one marking.  All local-system terms and unmultiplied pullback terms then
vanish in degrees strictly above the relevant virtual cohomological
dimensions, while the remaining term is a cotangent-line class multiplied by
lowest-weight cohomology at the inclusive endpoint $n-1=c(g)$.  This proves
that all three open inputs have the required Tate type.

The proof of Theorem~\ref{thm:range-intro} then follows the boundary induction
of Canning--Larson--Payne, but it is useful to give the complete argument.
On a product boundary atlas, an odd--odd K\"unneth summand in total degree
$16$ has one factor in odd degree at most seven and therefore vanishes by
Bergstr\"om--Faber--Payne \cite{BFP}.  A positive even--even summand has both
degrees at most fourteen and is Tate by \cite[Theorem~1.5(3)]{CLP}.  The two
axis summands are supplied by induction.  The range in
\eqref{eq:range-intro} is closed under passage to the vertex types of a
one-edge stable graph, so the induction terminates.

Finally, Theorem~\ref{thm:range-intro} shows that the degree-sixteen
cohomology of every component of the normalized boundary of
$\Mbar_{6,6}$ is Tate.  Hence the image of boundary homology is a quotient of
a direct sum of copies of $\Q(8)$.  The cokernel is the open group in
Theorem~\ref{thm:open-intro}.  This makes $H_{16}(\Mbar_{6,6})$ an extension
of Tate Hodge structures of the same weight; duality gives
Theorem~\ref{thm:main}.

The conclusion concerns Hodge type only.  A pure rational Hodge structure of
type $(8,8)$ need not be known to be generated by algebraic classes merely
from its Hodge type.  Accordingly, we do not claim that
$H^{16}(\Mbar_{6,6})$ is tautological, nor do we make an integral statement
or a statement about coarse moduli spaces.

\section{Hodge-theoretic conventions and the boundary sequence}

\subsection{Weights, twists, and duality}

The Tate Hodge structure $\Q(m)$ has weight $-2m$ and Hodge type
$(-m,-m)$.  Thus $\Q(-8)$ has weight $16$ and type $(8,8)$, while
$\Q(8)$ has weight $-16$ and type $(-8,-8)$.  We use the increasing weight
filtration $W_\bullet$.  The mixed Hodge theory used below is due to Deligne
\cite{DeligneHodgeII,DeligneHodgeIII}; cohomology and duality for
Deligne--Mumford stacks are discussed, for example, by Behrend
\cite{BehrendStacks}.

If $X$ is smooth of pure complex dimension $d$, Poincar\'e duality gives
\begin{equation}\label{eq:pd-smooth}
 H_i^{\BM}(X;\Q)\cong H^{2d-i}(X;\Q)(d)
\end{equation}
as mixed Hodge structures.  If $X$ is proper, ordinary homology is the Hodge
dual of ordinary cohomology:
\begin{equation}\label{eq:proper-dual}
  H_i(X;\Q)\cong H^i(X;\Q)^\vee.
\end{equation}
For a smooth proper Deligne--Mumford stack, $H^i$ is polarizable and pure of
weight $i$, so $H_i$ is polarizable and pure of weight $-i$.

We repeatedly use two elementary consequences.  First, every subobject and
every quotient of $\Q(m)^{\oplus N}$ in the category of rational mixed Hodge
structures is again a finite direct sum of copies of $\Q(m)$.  This follows
directly from strictness of the Hodge and weight filtrations.  Second, the
category of polarizable pure rational Hodge structures of a fixed weight is
semisimple.  Hence an extension of finite direct sums of a fixed Tate Hodge
structure by finite direct sums of the same Tate Hodge structure splits
noncanonically.

\subsection{The normalized boundary}

Write $\widetilde{\partial\M}_{g,n}$ for the normalization of the reduced
boundary $\Mbar_{g,n}\setminus\M_{g,n}$, and let
\[
  \nu\colon\widetilde{\partial\M}_{g,n}\longrightarrow\Mbar_{g,n}
\]
be its natural finite map.  Knudsen's clutching constructions
\cite{Knudsen} give a decomposition
\begin{equation}\label{eq:boundary-graphs}
 \widetilde{\partial\M}_{g,n}
 \cong
 \coprod_{\substack{[\Gamma]\text{ stable}\
                     \Gamma\text{ of type }(g,n)\\ |E(\Gamma)|=1}}
 \left[\Mbar_\Gamma/\Aut(\Gamma)\right],
 \qquad
 \Mbar_\Gamma:=\prod_{v\in V(\Gamma)}
                  \Mbar_{g_v,n_v}.
\end{equation}
This is also the description used in
\cite[Section~2.1]{CLP}.  One may see the normalization directly from the
local deformation space of a nodal curve: if $t_e$ is the smoothing
parameter of a node $e$, then the boundary is locally given by
$\prod_e t_e=0$, and its normalization records the choice of one node.
Partial normalization at that distinguished node gives the one-edge graph
in \eqref{eq:boundary-graphs}.

For a finite group $G$ acting on a complex Deligne--Mumford stack $Y$, the
Cartan--Leray spectral sequence and rational averaging give
\begin{equation}\label{eq:quotient-invariants}
  H^*([Y/G];\Q)\cong H^*(Y;\Q)^G.
\end{equation}
Indeed, $H^p(G,V)=0$ for $p>0$ over $\Q$.  In particular, taking graph
automorphism invariants cannot introduce a new Hodge type.

\subsection{The right-exact boundary sequence}

The localization sequence in Borel--Moore homology, together with strictness
for weights, yields the following fundamental input.

\begin{proposition}[Canning--Larson--Payne]\label{prop:boundary-sequence}
For every stable pair $(g,n)$ and every integer $i$, there is a right-exact
sequence of mixed Hodge structures
\begin{equation}\label{eq:boundary-sequence}
 H_i(\widetilde{\partial\M}_{g,n};\Q)
 \xrightarrow{\nu_*}
 H_i(\Mbar_{g,n};\Q)
 \longrightarrow
 W_{-i}H_i^{\BM}(\M_{g,n};\Q)
 \longrightarrow0.
\end{equation}
\end{proposition}

\begin{proof}
This is Equation~(6.1) of \cite{CLP}.  Notice that the first map is ordinary
proper pushforward in degree $i$.  There is no degree shift and no Tate twist
in \eqref{eq:boundary-sequence}.
\end{proof}

At $i=16$, let $I_{g,n}$ be the image of the first map.  Then
\eqref{eq:boundary-sequence} gives a short exact sequence
\begin{equation}\label{eq:boundary-short}
 0\longrightarrow I_{g,n}
 \longrightarrow H_{16}(\Mbar_{g,n};\Q)
 \longrightarrow W_{-16}H_{16}^{\BM}(\M_{g,n};\Q)
 \longrightarrow0.
\end{equation}
Thus the degree-sixteen problem separates into the Hodge type of a boundary
image and the Hodge type of one precise lowest-weight group on the open
moduli stack.

\section{The three endpoint calculations on the open moduli stacks}

\subsection{Forgetful and cotangent-line subspaces}

For $1\leq i\leq n$, let
\[
  \pi_i\colon\M_{g,n}\longrightarrow\M_{g,n-1}
\]
forget the $i$th marking, and let $\psi_i\in H^2(\M_{g,n};\Q)$ be the first
Chern class of the corresponding cotangent line.  Following
\cite[Sections~2 and 3]{CLP}, define
\begin{align}
 \Phi^k_{g,n}
   &:=\sum_{i=1}^n
       \pi_i^*W_kH^k(\M_{g,n-1};\Q),\label{eq:Phi-new}\\
 \Psi^k_{g,n}
   &:=\sum_{i=1}^n
       \psi_i\smile\Phi^{k-2}_{g,n}.
       \label{eq:Psi-new}
\end{align}
The first subspace consists of lowest-weight classes inherited from one
fewer marking.  The second is generated by multiplying such inherited
classes by one cotangent-line class.

Let $\mathbb V_\lambda$ denote the irreducible symplectic local system on
$\M_g$ indexed by a partition $\lambda$, and let
$V_{\lambda^{\mathsf T}}$ be the corresponding representation of the
symmetric group.  We use the following published description of the part not
accounted for by \eqref{eq:Phi-new} and \eqref{eq:Psi-new}.

\begin{proposition}[Canning--Larson--Payne]\label{prop:primitive-new}
If $g\geq2$, there is an isomorphism of pure rational Hodge structures
\begin{equation}\label{eq:primitive-new}
 \frac{W_kH^k(\M_{g,n};\Q)}
      {\Phi^k_{g,n}+\Psi^k_{g,n}}
 \cong
 \bigoplus_{|\lambda|=n}
 W_kH^{k-n}(\M_g;\mathbb V_\lambda)
 \otimes V_{\lambda^{\mathsf T}}.
\end{equation}
\end{proposition}

\begin{proof}
This is \cite[Lemma~3.1(a)]{CLP}.  Its proof uses the smooth-proper Leray
decomposition for powers of the closed universal curve.  In particular, the
identification of the index-two summands modulo the forgetful summands is
made using the relative Chow--K\"unneth projectors described by
Petersen--Tavakol--Yin \cite[Sections~5.1 and 5.2.2]{PTY}; it is not a
consequence merely of restricting a cotangent-line class to a fibre.
\end{proof}

\subsection{Virtual cohomological dimension with rational coefficients}

We record the coefficient version of the virtual cohomological dimension
vanishing that will be used in applying
Proposition~\ref{prop:primitive-new}.

\begin{lemma}\label{lem:vcd-rational}
Let $\Gamma$ be a virtually torsion-free group of virtual cohomological
dimension $d$, and let $L$ be a finite-dimensional rational
$\Gamma$-module.  Then
\[
  H^q(\Gamma;L)=0\qquad\text{for }q>d.
\]
\end{lemma}

\begin{proof}
Choose a torsion-free normal subgroup $\Gamma'\triangleleft\Gamma$ of finite
index and cohomological dimension $d$.  The group
$H^q(\Gamma';L)$ vanishes for $q>d$.  In the Hochschild--Serre spectral
sequence for
\[
  1\longrightarrow\Gamma'\longrightarrow\Gamma
   \longrightarrow F\longrightarrow1,
\]
the quotient $F$ is finite.  Taking $F$-invariants is exact over $\Q$, so
$H^p(F,U)=0$ for $p>0$ and every rational $F$-module $U$.  Consequently
\[
  H^q(\Gamma;L)\cong H^q(\Gamma';L)^F=0
  \qquad(q>d).
\]
See \cite[Chapter~VII]{Brown} for these standard group-cohomological facts.
\end{proof}

Teichm\"uller space is contractible, so rational cohomology of $\M_g$ and
$\M_{g,n}$, with the indicated local coefficients, agrees with the group
cohomology of the mapping class groups $\Mod_g$ and $\PMod_{g,n}$.
Harer computed, for $g\geq2$ and $n>0$,
\begin{equation}\label{eq:harer-vcd}
  \vcd(\Mod_g)=4g-5,
  \qquad
  \vcd(\PMod_{g,n})=4g-4+n
\end{equation}
\cite[Theorem~4.1]{Harer}.

\subsection{The endpoint proposition}

The numerical coincidences used in the proof are displayed in
Table~\ref{tab:endpoints}.  Here $d=3g-3+n$ and $c(g)$ is the bound in
\cite[Table~1]{CLP}.

\begin{table}[ht]
\centering
\caption{The three endpoint cases.}\label{tab:endpoints}
\begin{tabular}{@{}ccccc@{}}
\toprule
$g$ & $n$ & $d$ & $k=2d-16=4g-4+n$ & $c(g)=n-1$\\
\midrule
$3$ & $12$ & $18$ & $20$ & $11$\\
$5$ & $8$  & $20$ & $24$ & $7$\\
$6$ & $6$  & $21$ & $26$ & $5$\\
\bottomrule
\end{tabular}
\end{table}

\begin{proposition}\label{prop:endpoint-open}
For each triple $(g,n,k)$ in Table~\ref{tab:endpoints}, there is a
nonnegative integer $r_{g,n}$ such that
\begin{equation}\label{eq:endpoint-open}
  W_kH^k(\M_{g,n};\Q)
  =\RH^k(\M_{g,n};\Q)
  \cong\Q(-k/2)^{\oplus r_{g,n}}.
\end{equation}
\end{proposition}

\begin{proof}
Fix one of the three rows and apply
Proposition~\ref{prop:primitive-new}.  Since $k=4g-4+n$, the cohomological
degree on the right-hand side of \eqref{eq:primitive-new} is
\[
  k-n=4g-4>4g-5=\vcd(\Mod_g).
\]
Lemma~\ref{lem:vcd-rational} therefore makes every local-system summand
vanish.  Hence
\begin{equation}\label{eq:endpoint-phipsi}
  W_kH^k(\M_{g,n};\Q)=\Phi^k_{g,n}+\Psi^k_{g,n}.
\end{equation}

The unmultiplied pullback term also vanishes.  Indeed,
\[
 \vcd(\PMod_{g,n-1})
   =4g-4+(n-1)=k-1<k,
\]
so Lemma~\ref{lem:vcd-rational} gives
$H^k(\M_{g,n-1};\Q)=0$.  It follows from
\eqref{eq:Phi-new} that $\Phi^k_{g,n}=0$, and therefore
\begin{equation}\label{eq:endpoint-psi}
  W_kH^k(\M_{g,n};\Q)=\Psi^k_{g,n}.
\end{equation}

It remains to type the source of $\Psi^k_{g,n}$.  In every row,
$n-1=c(g)$.  The bound is inclusive.  By
\cite[Table~1, Lemma~4.3, and Proposition~4.5]{CLP},
\[
  W_{k-2}H^{k-2}(\M_{g,n-1};\Q)
   =\RH^{k-2}(\M_{g,n-1};\Q).
\]
This group is generated by codimension-$(k-2)/2$ algebraic cycle classes and
is therefore a finite direct sum of copies of $\Q(-(k-2)/2)$.  Each
$\psi_i$ is an algebraic class of type $\Q(-1)$.  Cup product is a morphism
of Hodge structures, so \eqref{eq:Psi-new} and
\eqref{eq:endpoint-psi} show that $W_kH^k(\M_{g,n})$ is a quotient of a
finite direct sum of copies of
\[
  \Q(-1)\otimes\Q(-(k-2)/2)=\Q(-k/2).
\]
Thus it is itself a finite direct sum of $\Q(-k/2)$.

The same argument shows that every class in the group is tautological, so
$W_kH^k\subseteq\RH^k$.  Conversely, a degree-$k$ tautological cycle class
has Hodge type $(k/2,k/2)$ and weight $k$, hence belongs to $W_kH^k$.
This proves the equality and the isomorphism in \eqref{eq:endpoint-open}.
All virtual-cohomological-dimension inequalities used above are strict; no
vanishing is inferred from $k=\vcd(\M_{g,n})$.
\end{proof}

We now translate Proposition~\ref{prop:endpoint-open} into the form required
by the boundary sequence.

\begin{corollary}\label{cor:endpoint-bm}
For each of $(g,n)=(3,12),(5,8),(6,6)$,
\begin{equation}\label{eq:endpoint-bm}
 W_{-16}H_{16}^{\BM}(\M_{g,n};\Q)
 \cong\Q(8)^{\oplus r_{g,n}}.
\end{equation}
In particular, Theorem~\ref{thm:open-intro} holds with
$r=r_{6,6}$.
\end{corollary}

\begin{proof}
For the three rows of Table~\ref{tab:endpoints}, we have $k=2d-16$.
Taking the lowest-weight piece in \eqref{eq:pd-smooth} gives
\[
 W_{-16}H_{16}^{\BM}(\M_{g,n};\Q)
 \cong W_kH^k(\M_{g,n};\Q)(d).
\]
By Proposition~\ref{prop:endpoint-open}, the right-hand side is a direct sum
of copies of $\Q(-k/2)(d)$.  Since $k/2=d-8$,
\[
  \Q(-k/2)(d)=\Q(-(d-8)+d)=\Q(8).
\]
Concretely, the three twist calculations are
\begin{align*}
 \Q(-10)(18)&=\Q(8),&
 \Q(-12)(20)&=\Q(8),&
 \Q(-13)(21)&=\Q(8).
\end{align*}
The final case, together with \eqref{eq:endpoint-open}, is exactly
Theorem~\ref{thm:open-intro}.
\end{proof}

\section{A Tate range in degree sixteen}

We next prove Theorem~\ref{thm:range-intro}.  This section is independent of
the compact $(6,6)$ conclusion: the pair $(6,6)$ itself does not lie in the
range considered below.

\subsection{Published low-degree inputs}

We will use the following three results about compact moduli stacks.

\begin{proposition}\label{prop:low-degree-inputs}
For all stable pairs $(g,n)$:
\begin{enumerate}[label=\textup{(\roman*)}]
\item $H^j(\Mbar_{g,n};\Q)=0$ for odd $j\leq9$;
\item if $0\leq j\leq14$ is even, then
  $H^j(\Mbar_{g,n};\Q)$ is a finite direct sum of $\Q(-j/2)$;
\item for every $n$, $H^{16}(\Mbar_{1,n};\Q)$ and
  $H^{16}(\Mbar_{2,n};\Q)$ are finite direct sums of $\Q(-8)$.
\end{enumerate}
\end{proposition}

\begin{proof}
Part~(i) is the odd-cohomology vanishing theorem of
Bergstr\"om--Faber--Payne \cite[Theorem~1.1]{BFP}.  Part~(ii) follows from
\cite[Theorem~1.5(3) and the discussion following it]{CLP}.  For genus one,
Petersen proves that even cohomology is generated by boundary-stratum cycle
classes \cite[Theorem~1.1]{PetersenGenusOne}; in degree sixteen these have
type $(8,8)$.  For genus two, part~(iii) follows from
\cite[Corollary~2.7]{CLP}, applied with $k=8$.  Purity of the compact
cohomology turns these type statements into the displayed Tate
decompositions.
\end{proof}

The direct source of the open pure-weight calculations below is
\cite[Table~1 and Proposition~4.5]{CLP}.  Canning--Larson established the
geometric results behind most entries of that table \cite{CanningLarson}.
Here the table and proposition of \cite{CLP} are the direct locators.

\subsection{The numerical range and its boundary vertices}

Let $\mathcal R$ be the set of stable pairs satisfying
\begin{equation}\label{eq:range}
  \mathcal R
  :=\{(g,n):g\leq5,\ 2g+n\leq18\}
     \cup\{(6,n):0\leq n\leq5\}.
\end{equation}
The following elementary identity is the reason this range is compatible
with boundary induction.

\begin{lemma}\label{lem:marking-budget}
Let $\Gamma$ be a connected stable graph of type $(g,n)$.  Then
\begin{equation}\label{eq:budget}
 \sum_{v\in V(\Gamma)}(2g_v-2+n_v)=2g-2+n.
\end{equation}
\end{lemma}

\begin{proof}
Every leg occurs at one vertex and every edge contributes two half-edges, so
$\sum_v n_v=n+2|E|$.  The arithmetic genus formula gives
$\sum_v g_v=g-b_1(\Gamma)$, where
$b_1(\Gamma)=|E|-|V|+1$.  Substitution gives
\begin{align*}
 \sum_v(2g_v-2+n_v)
 &=2(g-|E|+|V|-1)-2|V|+n+2|E|\\
 &=2g-2+n.
\end{align*}
\end{proof}

\begin{lemma}\label{lem:range-closed}
Let $(g,n)\in\mathcal R$, and let $\Gamma$ be a one-edge stable graph of
type $(g,n)$.  Every vertex type $(g_v,n_v)$ of $\Gamma$ belongs to
$\mathcal R$.  Moreover, every factor $\Mbar_{g_v,n_v}$ has complex
dimension strictly smaller than $\dim\Mbar_{g,n}$.
\end{lemma}

\begin{proof}
Suppose first that $g\leq5$ and $2g+n\leq18$.  Genus cannot increase on a
vertex, so $g_v\leq5$.  Stability makes every summand on the left of
\eqref{eq:budget} a positive integer.  Hence
\[
  2g_v-2+n_v\leq2g-2+n,
\]
and therefore $2g_v+n_v\leq2g+n\leq18$.  Thus every vertex lies in the first
part of \eqref{eq:range}.

Now suppose $g=6$ and $n\leq5$.  For the nonseparating graph the unique
vertex has type $(5,n+2)$ and satisfies
$2\cdot5+(n+2)=12+n\leq17$, so it lies in the first part of
\eqref{eq:range}.  For a separating graph, the two positive summands in
\eqref{eq:budget} add to $10+n\leq15$, so each vertex satisfies
$2g_v+n_v\leq16$.  A vertex of genus at most five lies in the first part of
the range; a vertex of genus six then has at most four markings and lies in
the second part.

Finally, partial normalization at one node lowers the total complex
dimension by one.  In the nonseparating case the single factor has dimension
$3(g-1)-3+(n+2)=3g-4+n$.  In the separating case the factor dimensions add
to
\[
 (3g_1-3+n_1)+(3g_2-3+n_2)=3g-4+n.
\]
Each individual factor therefore has dimension smaller than
$3g-3+n$.
\end{proof}

\subsection{The open term throughout the range}

\begin{lemma}\label{lem:range-open}
Let $(g,n)\in\mathcal R$ with $g\in\{0,3,4,5,6\}$.  Then
\[
 W_{-16}H_{16}^{\BM}(\M_{g,n};\Q)
 \cong\Q(8)^{\oplus a_{g,n}}
\]
for some $a_{g,n}\geq0$.
\end{lemma}

\begin{proof}
Put $d=3g-3+n$ and $j=2d-16$.  If $j<0$, the dual group in
\eqref{eq:pd-smooth} is zero, and there is nothing to prove.  Assume
$j\geq0$.

The marking bounds in \cite[Table~1]{CLP} are
\[
 c(0)=\infty,\qquad c(3)=c(4)=11,
 \qquad c(5)=7,\qquad c(6)=5.
\]
The inequalities defining $\mathcal R$ give
\[
 \begin{array}{c|ccccc}
 g&0&3&4&5&6\\ \hline
 \text{largest possible }n&18&12&10&8&5.
 \end{array}
\]
Thus $n\leq c(g)$ except possibly at $(3,12)$ and $(5,8)$.  In the
nonexceptional cases, \cite[Proposition~4.5]{CLP} gives
\[
 W_jH^j(\M_{g,n};\Q)=\RH^j(\M_{g,n};\Q)
 \cong\Q(-j/2)^{\oplus a_{g,n}}.
\]
In the two exceptional cases the same conclusion is
Proposition~\ref{prop:endpoint-open}.  Applying \eqref{eq:pd-smooth} and
using $j=2d-16$ gives
\[
 W_{-16}H_{16}^{\BM}(\M_{g,n};\Q)
 \cong\Q(-j/2)^{\oplus a_{g,n}}(d)
 =\Q(8)^{\oplus a_{g,n}},
\]
as required.
\end{proof}

\subsection{K\"unneth decomposition in total degree sixteen}

\begin{lemma}\label{lem:kunneth-sixteen}
Let $(g_1,n_1)$ and $(g_2,n_2)$ be stable pairs.  Suppose, for $i=1,2$,
that $H^{16}(\Mbar_{g_i,n_i};\Q)$ is a finite direct sum of $\Q(-8)$.
Then
\[
 H^{16}(\Mbar_{g_1,n_1}\times\Mbar_{g_2,n_2};\Q)
 \cong\Q(-8)^{\oplus N}
\]
for some $N\geq0$.
\end{lemma}

\begin{proof}
The K\"unneth formula is an isomorphism of Hodge structures:
\[
 H^{16}(\Mbar_{g_1,n_1}\times\Mbar_{g_2,n_2};\Q)
 \cong
 \bigoplus_{p+q=16}
 H^p(\Mbar_{g_1,n_1};\Q)
 \otimes H^q(\Mbar_{g_2,n_2};\Q).
\]
If $p$ and $q$ are odd, then $\min(p,q)\leq7$, so one factor vanishes by
Proposition~\ref{prop:low-degree-inputs}(i).  If $p$ and $q$ are both
positive and even, then $2\leq p,q\leq14$.  Part~(ii) of that proposition
types the tensor product as a finite direct sum of
\[
  \Q(-p/2)\otimes\Q(-q/2)=\Q(-8).
\]
The only remaining terms have $(p,q)=(16,0)$ or $(0,16)$.  Since each
moduli stack is connected, $H^0\cong\Q(0)$, and the hypotheses type these
axis terms as sums of $\Q(-8)$.
\end{proof}

\subsection{Induction on dimension}

\begin{proof}[Proof of Theorem~\ref{thm:range-intro}]
We prove the assertion simultaneously for all pairs in $\mathcal R$ by
induction on $d=3g-3+n$.  If $d\leq7$, then $16>2d$ and
$H^{16}(\Mbar_{g,n};\Q)=0$.  This is the base case.

Suppose $d\geq8$.  If $g=1$ or $g=2$, the conclusion is
Proposition~\ref{prop:low-degree-inputs}(iii), with no restriction on $n$.
We may therefore assume $g\in\{0,3,4,5,6\}$.  Lemma~\ref{lem:range-open}
gives
\begin{equation}\label{eq:range-open-term}
 W_{-16}H_{16}^{\BM}(\M_{g,n};\Q)
 \cong\Q(8)^{\oplus a_{g,n}}.
\end{equation}

Consider a component $[\Mbar_\Gamma/\Aut(\Gamma)]$ of the normalized
boundary.  If $\Gamma$ is nonseparating, its ordered atlas is a single
factor $\Mbar_{g-1,n+2}$.  By Lemma~\ref{lem:range-closed}, this factor lies
in $\mathcal R$ and has smaller dimension, so the induction hypothesis makes
its degree-sixteen cohomology a direct sum of $\Q(-8)$.

If $\Gamma$ is separating, its ordered atlas is a product of two moduli
stacks.  Both vertex types lie in $\mathcal R$ at smaller dimensions by
Lemma~\ref{lem:range-closed}.  The induction hypothesis applies to the
degree-sixteen cohomology of each factor, and
Lemma~\ref{lem:kunneth-sixteen} types the total degree-sixteen cohomology of
the product.  In either case, taking the finite graph-automorphism invariants
in \eqref{eq:quotient-invariants} preserves the type $\Q(-8)$.  Summing over
the finitely many one-edge graphs gives
\begin{equation}\label{eq:range-boundary-cohom}
 H^{16}(\widetilde{\partial\M}_{g,n};\Q)
 \cong\Q(-8)^{\oplus N_{g,n}}.
\end{equation}
By proper duality \eqref{eq:proper-dual},
\begin{equation}\label{eq:range-boundary-hom}
 H_{16}(\widetilde{\partial\M}_{g,n};\Q)
 \cong\Q(8)^{\oplus N_{g,n}}.
\end{equation}

Let $I_{g,n}$ be the boundary image in
\eqref{eq:boundary-short}.  It is a quotient of the group in
\eqref{eq:range-boundary-hom}, hence
$I_{g,n}\cong\Q(8)^{\oplus e_{g,n}}$ for some $e_{g,n}$.  Combining this
with \eqref{eq:range-open-term}, the exact sequence
\eqref{eq:boundary-short} expresses $H_{16}(\Mbar_{g,n};\Q)$ as an
extension of copies of $\Q(8)$ by copies of $\Q(8)$.  This homology group is
the dual of the polarizable pure weight-sixteen Hodge structure
$H^{16}(\Mbar_{g,n};\Q)$, so it is polarizable and pure of weight $-16$.
Semisimplicity gives
\[
 H_{16}(\Mbar_{g,n};\Q)
 \cong\Q(8)^{\oplus(a_{g,n}+e_{g,n})}.
\]
Dualizing proves
\[
 H^{16}(\Mbar_{g,n};\Q)
 \cong\Q(-8)^{\oplus b_{g,n}}.
\]
The dimension has decreased at every use of the induction hypothesis, so the
induction is well founded.  This proves
Theorem~\ref{thm:range-intro}.
\end{proof}

\section{The normalized boundary in the genus-six, six-pointed case}

We now specialize the graph description to the boundary of the stack in the
main theorem.  Its complex dimension is
\begin{equation}\label{eq:dimension-66}
 \dim_{\C}\Mbar_{6,6}=3\cdot6-3+6=21,
\end{equation}
so every component of the normalized boundary has dimension twenty.

There is one nonseparating graph.  Partial normalization at its loop lowers
the genus by one and replaces the node by two markings.  The two branches
have no preferred order, and therefore the corresponding component is
\begin{equation}\label{eq:nonsep-component}
 \left[\Mbar_{5,8}/S_2\right],
\end{equation}
where $S_2$ exchanges the last two markings.

For a separating one-edge graph, partial normalization produces two stable
pointed curves with vertex types $(g_1,n_1)$ and $(g_2,n_2)$ satisfying
\begin{equation}\label{eq:separating-numerics}
 g_1+g_2=6,
 \qquad n_1+n_2=8.
\end{equation}
The two extra markings are the branches of the distinguished node.  The
component is the finite quotient
\begin{equation}\label{eq:separating-component}
 \left[
   (\Mbar_{g_1,n_1}\times\Mbar_{g_2,n_2})/
   \Aut(\Gamma)
 \right].
\end{equation}
Equations \eqref{eq:nonsep-component} and
\eqref{eq:separating-component} exhaust the one-edge graphs.

\begin{proposition}\label{prop:boundary-66-tate}
There is a nonnegative integer $N$ such that
\begin{align}
 H^{16}(\widetilde{\partial\M}_{6,6};\Q)
   &\cong\Q(-8)^{\oplus N},\label{eq:boundary-66-cohom}\\
 H_{16}(\widetilde{\partial\M}_{6,6};\Q)
   &\cong\Q(8)^{\oplus N}.
   \label{eq:boundary-66-hom}
\end{align}
\end{proposition}

\begin{proof}
The pair $(5,8)$ satisfies $5\leq5$ and
$2\cdot5+8=18$, so Theorem~\ref{thm:range-intro} gives
\[
 H^{16}(\Mbar_{5,8};\Q)\cong\Q(-8)^{\oplus b_{5,8}}.
\]
Taking $S_2$-invariants therefore proves the required type for the
nonseparating component \eqref{eq:nonsep-component}.

Now consider a separating graph.  Put $m_i=2g_i+n_i$.  From
\eqref{eq:separating-numerics},
\[
  m_1+m_2=2(g_1+g_2)+(n_1+n_2)=20.
\]
Stability gives $m_i\geq3$, so $m_i\leq17$.  If $g_i\leq5$, then
$(g_i,n_i)$ lies in the first part of the range \eqref{eq:range}.  If
$g_i=6$, then $n_i=m_i-12\leq5$, so it lies in the second part.  Thus
Theorem~\ref{thm:range-intro} applies to both factors of the ordered atlas.
Lemma~\ref{lem:kunneth-sixteen} and rational graph-automorphism invariants
then give the type $\Q(-8)$ for every separating component.

The graph set is finite, so the direct sum over all components proves
\eqref{eq:boundary-66-cohom}.  Proper duality
\eqref{eq:proper-dual} gives \eqref{eq:boundary-66-hom}.
\end{proof}

The boundary map itself may have a nontrivial kernel.  Its rank is not
needed, because its source has only one Hodge type.

\begin{corollary}\label{cor:boundary-image}
For some integer $m$ with $0\leq m\leq N$,
\[
 \im\!\left(
  H_{16}(\widetilde{\partial\M}_{6,6};\Q)
  \longrightarrow H_{16}(\Mbar_{6,6};\Q)
 \right)
 \cong\Q(8)^{\oplus m}.
\]
\end{corollary}

\begin{proof}
The image is a quotient, in the category of mixed Hodge structures, of the
source in \eqref{eq:boundary-66-hom}.  Every quotient of a finite direct sum
of $\Q(8)$ is again a finite direct sum of $\Q(8)$.
\end{proof}

It is worth noting that no cohomology group of $\Mbar_{6,6}$ was used to
prove Corollary~\ref{cor:boundary-image}.  The image is controlled entirely
from the normalized boundary, so the argument does not feed the desired
compact conclusion back into itself.

\section{Proof of the main theorem}

\begin{proof}[Proof of Theorem~\ref{thm:main}]
Specialize \eqref{eq:boundary-short} to $(g,n)=(6,6)$.  By
Corollary~\ref{cor:boundary-image}, its boundary image is
$\Q(8)^{\oplus m}$.  By Corollary~\ref{cor:endpoint-bm}, its open quotient
is $\Q(8)^{\oplus r}$.  We therefore have a short exact sequence of mixed
Hodge structures
\begin{equation}\label{eq:final-short}
 0\longrightarrow\Q(8)^{\oplus m}
 \longrightarrow H_{16}(\Mbar_{6,6};\Q)
 \longrightarrow\Q(8)^{\oplus r}
 \longrightarrow0.
\end{equation}

The stack $\Mbar_{6,6}$ is smooth and proper, so the middle term is a
polarizable pure Hodge structure of weight $-16$.  Semisimplicity splits
\eqref{eq:final-short} noncanonically and gives
\[
 H_{16}(\Mbar_{6,6};\Q)
 \cong\Q(8)^{\oplus(m+r)}.
\]
By proper duality,
\[
 H^{16}(\Mbar_{6,6};\Q)
 \cong H_{16}(\Mbar_{6,6};\Q)^\vee
 \cong\Q(-8)^{\oplus(m+r)}.
\]
The integer $m+r$ is the dimension of this cohomology group, hence equals
$b_{16}$.  This proves the theorem.
\end{proof}

\section{Scope of the result}

\begin{remark}[Hodge type versus algebraic generation]
Theorem~\ref{thm:main} says that the complete degree-sixteen cohomology group
has Hodge type $(8,8)$.  The argument does not exhibit a basis of cycle
classes on $\Mbar_{6,6}$.  Thus it does not prove that
$H^{16}(\Mbar_{6,6})$ is tautological and does not, by Hodge type alone,
prove surjectivity of the cycle class map in codimension eight.
\end{remark}

\begin{remark}[Ranks]
The proof introduces three finite ranks: $r$ for the open quotient, $N$ for
the normalized boundary, and $m$ for the image of boundary homology.  It
shows $b_{16}=m+r$ and $0\leq m\leq N$, but computes none of these integers.
The theorem includes the possibility that any of the corresponding summands
is zero.
\end{remark}

\begin{remark}[Coefficients and stacks]
All statements are over $\C$ with rational coefficients and concern the
cohomology of Deligne--Mumford stacks.  Rational coefficients are used
essentially when passing to invariants of finite graph automorphism groups.
No integral assertion and no comparison with the cohomology of the coarse
moduli space is made.
\end{remark}

\begin{remark}[Why the endpoint is sharp for this proof]
The open calculation for $\M_{6,6}$ uses the equality $6-1=c(6)=5$; the
calculations for $\M_{5,8}$ and $\M_{3,12}$ similarly use
$8-1=c(5)=7$ and $12-1=c(3)=11$.  The argument reaches one marking beyond
the published Chow--K\"unneth range only because the primitive quotient and
the unmultiplied pullback term vanish strictly above virtual cohomological
dimension.  No Chow--K\"unneth generation statement is asserted for the
three endpoint stacks themselves.
\end{remark}

\section*{Acknowledgments}

The principal proof development, mathematical analysis, drafting, and review
underlying this manuscript were carried out by the artificial intelligence
collaborators named in the byline.

\begin{thebibliography}{99}

\bibitem{BehrendStacks}
K.~Behrend,
\emph{Cohomology of stacks},
in \emph{Intersection theory and moduli}, ICTP Lecture Notes, vol.~XIX,
Abdus Salam International Centre for Theoretical Physics, Trieste, 2004,
pp.~249--294.

\bibitem{BFP}
J.~Bergstr\"om, C.~Faber, and S.~Payne,
\emph{Polynomial point counts and odd cohomology vanishing on moduli spaces
of stable curves},
Ann. of Math. (2) \textbf{199} (2024), no.~3, 1323--1365,
\href{https://doi.org/10.4007/annals.2024.199.3.7}
{doi:10.4007/annals.2024.199.3.7}.

\bibitem{Brown}
K.~S. Brown,
\emph{Cohomology of groups},
Graduate Texts in Mathematics, vol.~87, Springer-Verlag, New York, 1982,
\href{https://doi.org/10.1007/978-1-4684-9327-6}
{doi:10.1007/978-1-4684-9327-6}.

\bibitem{CanningLarson}
S.~Canning and H.~Larson,
\emph{On the Chow and cohomology rings of moduli spaces of stable curves},
J. Eur. Math. Soc. \textbf{28} (2026), no.~9, 3869--3918,
\href{https://doi.org/10.4171/JEMS/1543}{doi:10.4171/JEMS/1543}.

\bibitem{CLP}
S.~Canning, H.~Larson, and S.~Payne,
\emph{Extensions of tautological rings and motivic structures in the
cohomology of $\overline{\mathcal M}_{g,n}$},
Forum Math. Pi \textbf{12} (2024), article no.~e23, 31 pp.,
\href{https://doi.org/10.1017/fmp.2024.24}{doi:10.1017/fmp.2024.24}.

\bibitem{DeligneHodgeII}
P.~Deligne,
\emph{Th\'eorie de Hodge. II},
Publ. Math. Inst. Hautes \'Etudes Sci. \textbf{40} (1971), 5--57,
\href{https://doi.org/10.1007/BF02684692}{doi:10.1007/BF02684692}.

\bibitem{DeligneHodgeIII}
P.~Deligne,
\emph{Th\'eorie de Hodge. III},
Publ. Math. Inst. Hautes \'Etudes Sci. \textbf{44} (1974), 5--77,
\href{https://doi.org/10.1007/BF02685881}{doi:10.1007/BF02685881}.

\bibitem{Harer}
J.~L. Harer,
\emph{The virtual cohomological dimension of the mapping class group of an
orientable surface},
Invent. Math. \textbf{84} (1986), no.~1, 157--176,
\href{https://doi.org/10.1007/BF01388737}{doi:10.1007/BF01388737}.

\bibitem{Knudsen}
F.~F. Knudsen,
\emph{The projectivity of the moduli space of stable curves. II. The stacks
$\overline{\mathcal M}_{g,n}$},
Math. Scand. \textbf{52} (1983), 161--199,
\href{https://doi.org/10.7146/math.scand.a-12001}
{doi:10.7146/math.scand.a-12001}.

\bibitem{PayneSurvey}
S.~Payne,
\emph{On the Hodge and Tate conjectures for moduli spaces of curves},
2026,
\href{https://arxiv.org/abs/2605.20453}{arXiv:2605.20453v1}.

\bibitem{PetersenGenusOne}
D.~Petersen,
\emph{The structure of the tautological ring in genus one},
Duke Math. J. \textbf{163} (2014), no.~4, 777--793,
\href{https://doi.org/10.1215/00127094-2429916}
{doi:10.1215/00127094-2429916}.

\bibitem{PTY}
D.~Petersen, M.~Tavakol, and Q.~Yin,
\emph{Tautological classes with twisted coefficients},
Ann. Sci. \'Ec. Norm. Sup\'er. (4) \textbf{54} (2021), no.~5, 1179--1236,
\href{https://doi.org/10.24033/asens.2479}{doi:10.24033/asens.2479}.

\end{thebibliography}

\end{document}
